\documentclass[aspectratio=169,11pt]{beamer}

% Idioma e codificação: compatíveis com pdfLaTeX.
\usepackage[T1]{fontenc}
\usepackage[utf8]{inputenc}
\usepackage[brazilian]{babel}
\usepackage{lmodern}
\usepackage{microtype}
\usepackage{csquotes}

% Matemática, figuras, tabelas e diagramas.
\usepackage{amsmath,amssymb,mathtools}
\usepackage{graphicx}
\usepackage{booktabs}
\usepackage{tikz}
\usetikzlibrary{arrows.meta,positioning,fit,calc,shapes.geometric}

% Bibliografia: BibLaTeX + Biber (não BibTeX).
\usepackage[
  backend=biber,
  style=authoryear,
  sorting=nyt,
  giveninits=true,
  uniquename=init,
  maxcitenames=2,
  maxbibnames=99,
  doi=true,
  url=true,
  isbn=false
]{biblatex}
\addbibresource{references.bib}

% Tema inicial deliberadamente simples; personalize após aprovar a narrativa.
\usetheme{Madrid}
\usecolortheme{dolphin}
\setbeamertemplate{navigation symbols}{}
\setbeamertemplate{caption}[numbered]
\setbeamertemplate{itemize items}[circle]
\setbeamertemplate{enumerate items}[default]
\setbeamersize{text margin left=8mm,text margin right=8mm}

% Metadados: substitua os campos antes da apresentação.
\title[Soberania digital e computação]{Soberania digital e computação}
\subtitle{Dependências, capacidades e autonomia tecnológica}
\author{[Nome do apresentador]}
\institute{[Instituição]}
\date{[Evento]\par [Data]}

% Links discretos; o próprio Beamer já carrega hyperref.
\hypersetup{
  colorlinks=true,
  linkcolor=blue,
  citecolor=blue,
  urlcolor=blue,
  pdftitle={Soberania digital e computação},
  pdfauthor={[Nome do apresentador]}
}

% Indique a duração estimada e a função do slide nas notas.
\newcommand{\slidepurpose}[1]{\note{\textbf{Função narrativa:} #1}}
\newcommand{\placeholder}[1]{\textit{\color{gray}#1}}

% Para gerar notas em páginas separadas, substitua hide notes por show notes.
\setbeameroption{hide notes}

\begin{document}

\begin{frame}[plain]
  \titlepage
  \note{Apresente a pergunta que orienta a palestra. Não antecipe toda a conclusão.}
\end{frame}

\section{Problematização}

\begin{frame}{De quais partes depende um serviço digital?}
  \begin{columns}[T,onlytextwidth]
    \column{0.56\textwidth}
      \begin{itemize}
        \item Aplicação e dados.
        \item Sistema operacional e bibliotecas.
        \item Rede, identidade e armazenamento.
        \item Infraestrutura física e operação.
        \item Contratos, competências e governança.
      \end{itemize}
    \column{0.40\textwidth}
      \begin{block}{Pergunta}
        Quem consegue compreender, operar, modificar e substituir cada componente?
      \end{block}
  \end{columns}
  \note{Use aqui um caso concreto que seja familiar ao público. O diagrama é apenas uma lista inicial de dimensões, não uma arquitetura universal.}
\end{frame}

\section{Conceitos}

\begin{frame}{Soberania digital: uma definição de trabalho}
  \begin{block}{Definição provisória}
    Capacidade técnica, institucional e política para compreender, operar, modificar e governar as infraestruturas digitais das quais se depende.
  \end{block}
  \begin{alertblock}{Distinção importante}
    Soberania digital não é sinônimo de isolamento tecnológico, segurança absoluta ou localização física de dados.
  \end{alertblock}
  \note{A definição é provisória e deve ser substituída ou qualificada de acordo com a literatura selecionada na especificação. Explique que soberania é um conceito disputado.}
\end{frame}

\section{Infraestrutura computacional}

\begin{frame}{As camadas de uma infraestrutura computacional}
  \centering
  \begin{tikzpicture}[
    layer/.style={draw,rounded corners,minimum width=8.5cm,minimum height=0.62cm,align=center},
    flow/.style={-{Latex[length=2mm]},thick}
  ]
    \node[layer,fill=blue!8] (users) {Usuários, serviços e aplicações};
    \node[layer,fill=blue!12,below=3mm of users] (runtime) {Bibliotecas, runtimes e sistemas operacionais};
    \node[layer,fill=blue!16,below=3mm of runtime] (infra) {Computação, armazenamento e virtualização};
    \node[layer,fill=blue!20,below=3mm of infra] (network) {Redes, conectividade e identidade};
    \node[layer,fill=blue!24,below=3mm of network] (physical) {Hardware, firmware e instalações};
    \draw[flow] (users) -- (runtime);
    \draw[flow] (runtime) -- (infra);
    \draw[flow] (infra) -- (network);
    \draw[flow] (network) -- (physical);
  \end{tikzpicture}
  \vspace{0.5em}
  \small\placeholder{Modelo analítico inicial; adaptar ao caso e às fontes escolhidas.}
  \note{A separação em camadas ajuda a localizar dependências, mas as relações reais são muitas vezes horizontais e circulares. O desenho não representa uma pilha universal.}
\end{frame}

\section{Dependências}

\begin{frame}{Dependências não são apenas técnicas}
  \begin{columns}[T,onlytextwidth]
    \column{0.48\textwidth}
      \begin{block}{Dimensões}
        \begin{itemize}
          \item Técnica
          \item Econômica
          \item Jurídica
          \item Operacional
          \item Conhecimento e competências
        \end{itemize}
      \end{block}
    \column{0.48\textwidth}
      \begin{block}{Questões de análise}
        \begin{itemize}
          \item Quem controla?
          \item Quem pode auditar?
          \item Quem consegue manter?
          \item Quão difícil é migrar?
          \item Há alternativas interoperáveis?
        \end{itemize}
      \end{block}
  \end{columns}
  \note{Substitua ou complemente esta lista com evidências de estudos de caso. Diferencie a existência de uma alternativa da capacidade real de migrar para ela.}
\end{frame}

\section{Capacidades}

\begin{frame}{Controle exige capacidades concretas}
  \begin{enumerate}
    \item Conhecer e documentar a infraestrutura.
    \item Auditar componentes e cadeias de fornecimento.
    \item Operar e manter os serviços.
    \item Corrigir, modificar e adaptar sistemas.
    \item Migrar e substituir componentes.
    \item Definir regras de governança e responsabilização.
  \end{enumerate}
  \note{Destaque que possuir o código-fonte ou os dados não significa, automaticamente, ter equipe, orçamento, documentação ou conhecimento para operar a infraestrutura.}
\end{frame}

\section{Alternativas e limites}

\begin{frame}{Software livre e padrões abertos podem ampliar opções}
  \begin{block}{Possíveis contribuições}
    \begin{itemize}
      \item Estudo, auditoria e modificação do software.
      \item Redistribuição e adaptação conforme a licença \parencite{stallman2002free}.
      \item Interoperabilidade e maior portabilidade, quando efetivamente implementadas.
    \end{itemize}
  \end{block}
  \begin{alertblock}{O que não garantem isoladamente}
    Capacidade operacional, independência de hardware, ausência de custos, segurança automática ou governança adequada.
  \end{alertblock}
  \note{Use fontes primárias para definir software livre e fontes técnicas para discutir padrões. Evite apresentar software livre como solução suficiente para todos os problemas de soberania.}
\end{frame}

\section{Síntese}

\begin{frame}{Perguntas para avaliar uma infraestrutura}
  \begin{enumerate}
    \item Sabemos de quais componentes e fornecedores dependemos?
    \item Conseguimos auditar e operar os componentes críticos?
    \item Temos competências para corrigir e manter os sistemas?
    \item Existem alternativas e caminhos de migração testados?
    \item Quem define as regras e responde por decisões e incidentes?
  \end{enumerate}
  \note{Retome o exemplo da abertura e aplique estas perguntas. Deixe claro que se trata de uma heurística de análise, não de uma certificação formal.}
\end{frame}

\begin{frame}{Conclusões}
  \begin{itemize}
    \item Soberania digital é multidimensional.
    \item Dependências atravessam camadas técnicas e institucionais.
    \item Autonomia requer capacidades concretas, alternativas e governança.
  \end{itemize}
  \vspace{1em}
  \begin{block}{Pergunta para discussão}
    Quais componentes críticos da infraestrutura que utilizamos conseguimos efetivamente operar e substituir?
  \end{block}
  \note{Ajuste estas conclusões ao conjunto de evidências que de fato aparecerá na versão final. Não apresente como resultados demonstrados aquilo que for apenas hipótese.}
\end{frame}


\begin{frame}[allowframebreaks]{Referências}
  \small
  \printbibliography[heading=none]
\end{frame}

\begin{frame}[plain]
  \centering
  {\LARGE Obrigado!}\par\vspace{1em}
  {\large Perguntas e discussão}\par\vspace{2em}
  \placeholder{[Nome, instituição e contato]}
\end{frame}

\end{document}
